\subsection{实验原理}
任务二的核心对象是向量求和程序，其基本思想是对向量中的各个元素依次进行读取并累加，最终得到整个向量的和。程序正确性的判断依据是：在给定固定输入数据时，优化前后各版本都应输出相同的求和结果。
\begin{itemize}
    \item 程序性能分析的基本原理:本实验不仅要求程序能够正确运行，还要求对程序执行效率进行分析。实验中通过观察程序运行过程中的周期数等性能指标，比较不同实现方式下程序的执行开销，从而评价优化方法是否有效。
    \item 减少过程调用的优化原理:在优化过程中，可以将原先通过函数调用完成的操作改为直接访问结构体成员和数组元素，例如直接使用 $v->len$ 和 $v->data[i]$。这样可以减少循环中的过程调用次数，缩短执行路径，降低不必要的开销。
    \item 局部累加变量的优化原理:若在循环中每次都通过 $*dest = *dest + val$ 更新结果，会引入额外的存储访问。通过设置局部变量 sum 或 t 在寄存器中完成累加，最后再一次性写回目标地址，可以减少循环体中的存储操作，提高运行效率。
    \item 循环展开的优化原理:循环展开的思想是每次循环不再只处理一个元素，而是同时处理多个元素。例如 4 路展开或 8 路展开，可以减少循环判断、分支跳转和计数更新等控制开销，使更多指令用于有效计算。
    \item 多累加器的优化原理:在循环展开基础上，使用多个独立累加变量，可以减轻单一累加变量带来的数据相关问题，降低顺序依赖对流水线执行效率的限制，从而进一步提升程序性能。这也是优化二和优化四取得更好效果的重要原因。
\end{itemize}